\chapter{Physics and algorithms}
\label{ch:physics}

This chapter explains the calculations that all the analysis modules share. Chapter~\ref{ch:geometry} defines the geometry.

\section{Attenuation: the Beer--Lambert law}

A particle that travels a distance $L_i$ in a material with linear attenuation coefficient $\mu_i$ has the probability
$\exp(-\mu_iL_i)$ of crossing it without interacting. After crossing several layers,
\begin{equation}
  T = \exp\!\Big(-\sum_i \mu_i L_i\Big) = \e^{-\tau},\qquad \tau=\sum_i\mu_iL_i,
  \label{eq:beer}
\end{equation}
where $\tau$ is the \emph{optical depth}. This is what \code{sample.transmission(path_lengths)} computes, from a table of
path lengths with one row per particle and one column per layer.

\section{Where does a particle interact?}

The optical depth at which a particle has its first interaction is a random number with the exponential distribution. It is
generated from a uniform random number $u\in[0,1)$ with
\begin{equation}
  \tau = -\ln(1-u).
  \label{eq:tau}
\end{equation}
The particle interacts inside the sample if $\tau$ is smaller than the total optical depth of its path through the sample;
otherwise it crosses the sample without interacting. To convert $\tau$ into a position one follows the particle through the
layers, adding $\mu_iL_i$ until the value $\tau$ is reached; inside that layer the position is found by linear interpolation
(the optical depth grows linearly with distance, with slope $\mu_i$).

\subsection{Slab}

For a slab rotated by $\omega=90^\circ-\alpha$ (Chapter~\ref{ch:geometry}), the beam is oblique with respect to the layers:
the path through a thickness $t$ is $t/\cos\omega$, and the footprint of the beam on the front face is stretched by
$1/\cos\omega$. The optical depth along the beam equals the optical depth along the normal divided by $\cos\omega$, so a random
$\tau$ corresponds to the value $\tau\cos\omega$ measured along the normal. The beam is a ray
\begin{equation}
  x' = x_\mathrm{entry}-\sigma\sin\omega,\qquad z' = z_\mathrm{front}+\sigma\cos\omega,
\end{equation}
as a function of the distance $\sigma$ along it. Because the section of the sample is finite, the function also
takes into account that a particle can \emph{enter through a side face}, and that it can \emph{leave through a side face before
interacting}: such particles are handled correctly (the optical depth only starts to accumulate at the entry point).
In practice the calculation follows each beam particle from a point upstream of the sample and finds, for every layer, the
interval of the path spent in its material; the same intervals serve for partially filled layers (Section~\ref{sec:partial-algorithm}).

\subsection{Disc}

The disc (Chapter~\ref{ch:disc}) is handled exactly as the slab, with one change: the box of the sample is replaced by the
cylinder $x'^2+y'^2\le R^2$ whose axis is the normal $z'$, between the front and back faces. Its intersection with a straight
line is one interval, the solution of a quadratic equation like the one of the next subsection, written in the plane $(x',y')$
of the disc instead of the horizontal plane $(x,z)$ (Section~\ref{sec:disc-algorithm}).

\subsection{Cylinder}

Seen from above, a cylinder of radius $R$ is the circle $x^2+z^2=R^2$. A particle that moves as $\vec p+t\,\hat d$ meets it
where
\begin{equation}
  a\,t^2+2b\,t+c=0,\quad a=d_x^2+d_z^2,\;\; b=x\,d_x+z\,d_z,\;\; c=x^2+z^2-R^2,
\end{equation}
that is, at $t=\big(-b\pm\sqrt{b^2-ac}\big)/a$. Applying this to the radius of every layer, and adding the places where the particle
reaches the top or bottom faces ($y=\pm h/2$), cuts the path into \emph{segments}. Inside a segment the material is always the
same, and it is found from the distance of the segment's midpoint to the axis. The length of the path in each layer is the sum of the
lengths of its segments, and the next interaction is found by accumulating $\mu\,\Delta t$ segment by segment until the random
$\tau$ is reached. Particles parallel to the axis (no horizontal motion) are handled as a special case. When some layers are
partially filled, the crossings of their levels are added to the segments and the segments above the level of their layer get
$\mu=0$ (Section~\ref{sec:partial-algorithm}).

Figure~\ref{fig:traj} (Chapter~\ref{ch:geometry}) illustrates these steps: the incident particles end at their first interaction point, and from there the path towards the detector is followed.

\subsection{Sphere}

The sphere (Chapter~\ref{ch:sphere}) is cut into segments exactly as the cylinder, with spheres instead of cylinders: the path
$\vec p+t\,\hat d$ meets the sphere of radius $r_i$ where $t^2 + 2(\vec p\cdot\hat d)\,t + |\vec p|^2 - r_i^2 = 0$, the layer of a
segment is given by the distance of its midpoint to the centre, and there are no flat ends (Section~\ref{sec:sphere-algorithm}).

\section{Path from a point to the exit of the sample}

Given an interaction point $\vec p$ and a direction $\hat d$ (for instance that of the detector), the function
\code{outgoing_path_lengths} returns the distance travelled in each layer until the particle leaves the sample.
For the \textbf{slab}, the distance to the exit is $t_\mathrm{exit}=\min_k (b_k-p_k)/d_k$, where $b_k$ is the face reached in each coordinate
(the one in the sense of $d_k$); the layer $i$, between the planes $z_i^\mathrm{lo}$ and $z_i^\mathrm{hi}$, is crossed between
$t_a=(z_i^\mathrm{lo}-p_z)/d_z$ and $t_b=(z_i^\mathrm{hi}-p_z)/d_z$, and the path inside it is the overlap of that interval with
$[0,t_\mathrm{exit}]$. For the \textbf{disc} the exit is where the path leaves the cylinder of radius $R$ or one of the faces,
whichever comes first. For the \textbf{cylinder} and the \textbf{sphere} the segments described above are used. Particles travelling parallel to the layers (slab) stay in the layer where they start.

\section{Scattering, absorption and multiple scattering}

At every interaction the particle is scattered with the probability \code{scatter_fraction} of the layer; otherwise it is
absorbed and lost. A scattered particle gets a new direction: the deflection angle $\chi$ relative to the previous direction
is drawn from an angular distribution, and the azimuth around the previous direction is uniform. By default
the distribution is isotropic ($\cos\chi$ uniform in $[-1,1]$). A Henyey--Greenstein distribution \cite{hg} is available,
\begin{equation}
  \cos\chi=\frac{1}{2g}\left[1+g^2-\left(\frac{1-g^2}{1-g+2gu}\right)^{2}\right],
\end{equation}
with $g>0$ forward peaked, $g<0$ backward peaked and $g=0$ isotropic. Any function \code{(n, rng)} $\to$ $n$ values of $\cos\chi$
can be used. After scattering, the particle flies on and either interacts again or leaves the sample
(Chapter~\ref{ch:scattering}).

\section{Example}

\pyexample{ex03_interaction_paths}

The first lines show that 70\,\% of the particles of the beam interact; most of the interactions are in the inner Fe cylinder, which is the most
attenuating one. The paths are averages over the interacting particles: on the way out towards $\twotheta=60^\circ$ they cross
on average 2.2~mm of Fe, 1.2~mm of Al and 0.7~mm of Cu.
